% =====================================================================
%  2bat-ud5-9.tex  —  SESSIÓ 9: La distribució normal
%  (sense preàmbul: s'inclou des de main.tex)
% =====================================================================
\horatitol{Sessió 9 — La distribució normal}%
{Presentar la corba normal, els seus paràmetres (mitjana i desviació) i la idea que la probabilitat es llegeix com a àrea sota la corba.}

\subsection{Idea clau}

Ja hem vist que la binomial, amb moltes repeticions, dibuixa una \textbf{campana}. Avui
presentem el model continu que la descriu: la \textbf{distribució normal}. És la corba
més important de l'estadística perquè moltíssimes magnituds reals s'hi ajusten.

\begin{idea}
\textbf{La corba normal.} És una corba en forma de campana, \textbf{simètrica},
determinada per dos paràmetres:
\begin{itemize}[leftmargin=1.4em, itemsep=2pt]
  \item La \textbf{mitjana} $\mu$: marca el \textbf{centre} de la campana (on és el
        màxim).
  \item La \textbf{desviació típica} $\sigma$: mesura com d'\textbf{ampla} és la campana
        (com de repartides estan les dades). Poca $\sigma$ $\to$ campana estreta i alta;
        molta $\sigma$ $\to$ campana ampla i baixa.
\end{itemize}
S'escriu $X\sim N(\mu,\sigma)$. \textbf{Idea central:} la \textbf{probabilitat} d'un
interval és l'\textbf{àrea} sota la corba en aquell interval; l'àrea total val $1$
($100\%$).
\end{idea}

\begin{center}
\begin{tikzpicture}
\begin{axis}[udaxis, width=11cm, height=5cm,
    xlabel={\small $x$}, ylabel={}, ytick=\empty,
    xmin=-4, xmax=4, ymin=0, ymax=0.45,
    xtick={-3,-2,-1,0,1,2,3}, xticklabels={$\mu\!-\!3\sigma$,$\mu\!-\!2\sigma$,$\mu\!-\!\sigma$,$\mu$,$\mu\!+\!\sigma$,$\mu\!+\!2\sigma$,$\mu\!+\!3\sigma$}]
  \addplot[udblue, thick, domain=-4:4, samples=140]{0.3989*exp(-x^2/2)};
  \addplot[udblue!18, domain=-1:1, samples=80] {0.3989*exp(-x^2/2)} \closedcycle;
  \node[udnavy, font=\scriptsize] at (axis cs:0,0.16){$\approx 68\%$};
\end{axis}
\end{tikzpicture}

\figcap{La corba normal, centrada a $\mu$. L'àrea entre $\mu-\sigma$ i $\mu+\sigma$ és aproximadament el $68\%$ del total.}
\end{center}

\subsection{Exemples}

\begin{exemple}[llegir els paràmetres]
Les puntuacions d'una prova d'aptitud es reparteixen segons $N(50,10)$: la mitjana és
$\mu=50$ (el centre de la campana) i la desviació típica és $\sigma=10$ (l'amplada). La
majoria de puntuacions cauran a prop de $50$, i com més ens allunyem del centre, menys
freqüents són.
\end{exemple}

\begin{exemple}[la regla del 68 -- 95 -- 99{,}7]
En qualsevol distribució normal, sempre passa el mateix repartiment al voltant de la
mitjana:
\begin{itemize}[leftmargin=1.4em, itemsep=2pt]
  \item Entre $\mu-\sigma$ i $\mu+\sigma$ hi ha aproximadament el \textbf{68\%} de les
        dades.
  \item Entre $\mu-2\sigma$ i $\mu+2\sigma$, aproximadament el \textbf{95\%}.
  \item Entre $\mu-3\sigma$ i $\mu+3\sigma$, aproximadament el \textbf{99,7\%}.
\end{itemize}
Amb $N(50,10)$: el $68\%$ de les puntuacions estan entre $40$ i $60$; el $95\%$, entre
$30$ i $70$.
\end{exemple}

\subsection{Exercicis resolts}

\begin{exresolt}[aplicar la regla del 68]
L'alçada d'un grup segueix $N(170,8)$ (en cm). Entre quins valors hi ha aproximadament el
$68\%$ de les persones?
\solucio
El $68\%$ central està entre $\mu-\sigma$ i $\mu+\sigma$:
\[
170-8=162 \quad\text{i}\quad 170+8=178.
\]
Aproximadament el $68\%$ de les persones fan entre $162$ i $178$ cm.
\resposta{Entre $162$ i $178$ cm.}
\end{exresolt}

\begin{exresolt}[l'efecte de la desviació]
Dues proves tenen mitjana $\mu=50$, però una té $\sigma=5$ i l'altra $\sigma=15$. Quina
campana és més estreta? Què vol dir en termes de les dades?
\solucio
La de $\sigma=5$ és més \textbf{estreta i alta}: les dades estan més concentrades a prop
de $50$. La de $\sigma=15$ és més \textbf{ampla i baixa}: les dades estan més disperses.
Menys desviació vol dir més homogeneïtat.
\resposta{La de $\sigma=5$; indica dades més concentrades al voltant de la mitjana.}
\end{exresolt}

\subsection{Exercicis per fer}

\begin{exercicis}
\begin{exlist}
  \item En una $N(60,12)$, digues quant valen la mitjana i la desviació típica, i entre
        quins valors hi ha el $68\%$ central.
  \item En una $N(100,15)$, entre quins valors hi ha aproximadament el $95\%$ de les
        dades?
  \item Dues distribucions normals tenen la mateixa mitjana però desviacions $\sigma=4$ i
        $\sigma=10$. Quina és més ampla? Quina té les dades més disperses?
  \item L'àrea total sota una corba normal, quant val? Com es relaciona amb la
        probabilitat?
  \item \textbf{(Aplicació)} El temps d'atenció (en minuts) a un servei segueix
        $N(20,4)$. Entre quins valors espera estar el $68\%$ de les atencions? I el
        $95\%$?
  \item \textbf{(Interpretació)} Per què diem que la distribució normal és tan important?
        Posa dos exemples de magnituds reals que s'hi ajustin.
\end{exlist}
\end{exercicis}

% ---------------------------------------------------------------------
\fullsolucions{Sessió 9}
\begin{solucions}
\begin{sollist}
  \item $\mu=60$, $\sigma=12$. El $68\%$ central: entre $60-12=48$ i $60+12=72$.
  \item El $95\%$ està entre $\mu-2\sigma$ i $\mu+2\sigma$: $100-30=70$ i $100+30=130$.
        Entre $70$ i $130$.
  \item La de $\sigma=10$ és més ampla i té les dades més disperses; la de $\sigma=4$ és
        més estreta (dades més concentrades).
  \item L'àrea total val $1$ ($100\%$). La probabilitat d'un interval és l'àrea sota la
        corba en aquell interval, per això l'àrea de tot el conjunt és la probabilitat
        total, $1$.
  \item El $68\%$: entre $20-4=16$ i $20+4=24$ minuts. El $95\%$: entre $20-8=12$ i
        $20+8=28$ minuts.
  \item Resposta oberta. Idea: és important perquè moltíssimes magnituds reals s'hi
        ajusten i perquè permet calcular probabilitats de manera estàndard. Exemples:
        alçades, pesos, puntuacions de proves, errors de mesura, temps d'atenció\ldots
\end{sollist}
\end{solucions}
